\section{Introducción}
\label{sec:introduccion}

Cuando un grupo de personas trabaja de manera coordinada durante el tiempo suficiente,
recuerda más de lo que cualquiera de sus miembros podría
recordar de forma individual.
Wegner \cite{wegner1987transactive} llamó a este fenómeno
\textbf{memoria transactiva}: una forma de memoria compartida en la que cada individuo
almacena no solo sus propios conocimientos, sino también un \emph{directorio} que indica quién es reconocido
dentro del grupo como referente en cada dominio.
Cuando un miembro necesita información que no conoce, consulta ese directorio,
determina a quién preguntar y delega la recuperación.
El sistema funciona gracias a esa especialización: el conocimiento no necesita repetirse en cada miembro,
sino que se reparte entre ellos, y la coordinación entre pares reemplaza gradualmente a cualquier mediador central.

El modelo de Wegner identifica tres procesos principales.
Primero, la \emph{actualización del directorio}: tras cada interacción, todos los
miembros, incluso quienes no respondieron directamente, registran quién supo responder.
Segundo, la \emph{asignación de información}: la información nueva se envía
hacia el miembro reconocido como referencia en ese dominio.
Tercero, la \emph{coordinación de recuperación}: las consultas futuras viajan directamente
a ese referente sin pasar por un intermediario.
La transición de una \emph{fase temprana} mediada a una \emph{fase madura} punto a punto
es la señal de que el sistema transactivo ya quedó formado.

Un reporte previo de este proyecto mostró que un sistema de tres agentes especialistas
sobre memorias asociativas puras desarrolla el directorio transactivo de Wegner por
interacción acumulativa, alcanzando una precisión madura del 98.75\% en tres categorías
de ETH-80.
Este reporte técnico aborda la pregunta que ese resultado deja abierta:

\begin{quote}
¿La maquinaria transactiva construida sobre memorias asociativas entrópicas
\emph{escala} cuando el número de agentes especialistas crece de tres a ocho,
cubriendo todas las categorías de ETH-80? Y si aparecen límites, ¿provienen del
mecanismo asociativo o de las representaciones que lo alimentan?
\end{quote}

El sistema construido, denominado \textbf{EAM-TMS} (Sistema de Memoria Transactiva
Multi-Agente sobre Memorias Asociativas Entrópicas), se compone en esta versión de
\textbf{ocho agentes especialistas}, uno por cada categoría del conjunto ETH-80
\cite{leibe2003eth80}: \textit{apple}, \textit{car}, \textit{cow}, \textit{cup},
\textit{dog}, \textit{horse}, \textit{pear} y \textit{tomato}.
Un mediador central denominado TME (\textit{Transactive Memory Expert}) coordina la fase
temprana y registra en sus propios directorios qué agente responde a cada tipo de
consulta.
Conforme cada agente acumula registros en sus memorias de directorio internas
(una textual y una visual), el TME deja de hacer falta: en la fase madura, cualquier
agente puede determinar a quién reenviar una consulta desconocida, sea una frase o una
imagen, sin intervención del mediador.

El sistema no emplea ningún mecanismo de optimización numérica en la toma de decisiones.
Todas las operaciones (registro, reconocimiento, recuperación y enrutamiento) son
operaciones definidas formalmente sobre las \textbf{Memorias Asociativas Entrópicas}
de Pineda \& Morales en su variante ponderada (W-EAM)
\cite{pineda2021entropic,pineda2022weighted} y su extensión hetero-asociativa (EHAM)
\cite{morales2024hetero,morales2025missing}.
Las representaciones semánticas se obtienen con embeddings fastText
\cite{bojanowski2017fasttext}, cuantizados \emph{por magnitud} en 16 niveles
(un cambio central de esta versión frente a la binarización por signo anterior).
El vocabulario de llenado se extrae de ConceptNet 5.7 \cite{speer2017conceptnet}
y el texto se procesa con spaCy \cite{honnibal2020spacy}.
Las representaciones visuales las produce un encoder ResNet18 \cite{he2016resnet}
preentrenado, cuya salida se cuantiza y almacena.

\subsection*{Contribuciones principales}

Este trabajo aporta cinco contribuciones concretas respecto a la versión anterior:

\begin{enumerate}
    \item \textbf{Escalado a ocho agentes especialistas.}
    El sistema pasa de 3 a 8 categorías (las ocho clases de ETH-80), con el encoder
    visual reentrenado de forma determinista, el vocabulario semántico regenerado y el
    banco de evaluación extendido a 411 consultas ($\approx$50 por dominio).
    Los resultados muestran que los directorios, la lectura normalizada B1, el rechazo
    por containment y la coordinación punto a punto escalan sin cambios en el mecanismo.

    \item \textbf{Cuantización por magnitud de las pistas semánticas.}
    La binarización por signo de fastText (que usaba 2 de los 16 niveles disponibles)
    se reemplaza por una cuantización por magnitud $v/S$ sobre una escala global
    persistida, que puebla los 16 niveles.
    Este cambio de representación elevó la precisión temprana de 8 clases del 50\% al
    rango de 85--88\% sin tocar la maquinaria asociativa.

    \item \textbf{Directorio visual por agente.}
    Cada agente incorpora un directorio hetero-asociativo visual propio
    ($M^R_{\mathrm{dir}}$), de modo que una imagen puede entrar por \emph{cualquier}
    agente y ser redirigida punto a punto al especialista, replicando para la modalidad
    visual la coordinación de recuperación de Wegner.
    La alternativa de un puente por el TME traducía 0\% y fue descartada.

    \item \textbf{Auditoría de fidelidad del pipeline.}
    Se auditó que ninguna decisión provenga de mecanismos ajenos a las memorias: se
    eliminó el único vector sintético del vocabulario (que además distorsionaba la
    escala global de cuantización), se verificó que el código de Pineda \& Morales
    permanece byte-idéntico, y todos los números reportados provienen de corridas
    posteriores a esa auditoría, trazables a datos no fabricados.

    \item \textbf{Evaluación experimental del sistema escalado.}
    Se reportan el estudio de calibración del directorio (9 condiciones $\times$
    $N \in \{50,100,200,400\}$ $\times$ 5 semillas), el barrido de parámetros nativos
    $\iota \times \kappa$, la curva de formación del directorio, la curva de capacidad
    de llenado y el hemisferio visual completo, junto con el costo computacional de la
    reproducción íntegra.
\end{enumerate}

\subsection*{Estructura del documento}

El documento se organiza de la siguiente manera.
La Sección~\ref{sec:marco_teorico} presenta los fundamentos teóricos: la EAM de Pineda
\& Morales y su variante ponderada, la memoria hetero-asociativa 4D y el modelo de
memoria transactiva de Wegner.
La Sección~\ref{sec:arquitectura} describe en detalle la arquitectura EAM-TMS de ocho
agentes, incluyendo los diagramas de flujo del llenado, la fase temprana y la fase
madura, así como la formalización de todas las operaciones del sistema.
La Sección~\ref{sec:metodologia} detalla la metodología experimental.
La Sección~\ref{sec:resultados} reporta los resultados.
La Sección~\ref{sec:discusion} discute los resultados.
La Sección~\ref{sec:conclusiones} concluye y señala trabajo futuro.
